\section{Financial reporting quality}

\begin{qfmethod}[Reporting quality and economic results are separate dimensions]\label{meth:reporting-results-quality-spectrum}
\qffield{Statement} Reporting quality concerns whether information faithfully and usefully describes the business. Results quality concerns whether the reported earnings, cash generation and financial position are sustainable and adequate for the risk and capital employed. A truthful loss can be high-quality reporting of poor results; accurately disclosed disposal gains can be low-persistence earnings without fraudulent reporting. Conversely, apparently strong profits can be unusable if recognition is biased \citep[chap.~11, solutions~1--12, PDF pp.~217--218]{statement-analysis-workbook-2020}.
\qffield{Practice} Place reports on a continuum: useful compliant reporting of sustainable returns; compliant reporting of weak or transitory results; biased choices within permitted judgment; aggressive or inconsistent presentation that obscures economics; noncompliant recognition; and fictitious transactions. This ordering is a screening framework, not a numerical rating or proof of intent. Meridian's base profit of 112.5 accompanied by CFO 137.5 allows analysis, but the durability of its 17 percent operating margin remains a business question. In the later reporting variant, premature sales and capitalized routine maintenance overstate profit and impair reporting quality even though CFO still exceeds net income.
\qffield{Interpretation} Conservative estimates can depress today's results and create future releases; systematic conservatism is not the same as neutral reporting. Evaluate earnings quality only after establishing the accounting basis and information reliability. Use separate conclusions for reporting, earnings persistence, cash sustainability and balance-sheet recovery, so a clean audit opinion does not automatically become a favorable investment conclusion.
\end{qfmethod}

\begin{qfmethod}[Reporting incentives, opportunities and disciplinary mechanisms]\label{meth:reporting-incentives-controls}
\qffield{Statement} Bonuses, financing needs, covenant pressure, market expectations and job security can reward upward earnings management; new managers may instead take a large write-off to improve later comparisons. Opportunity comes from weak controls, complex transactions, estimation uncertainty and poor oversight, while rationalization can turn a target into a justification. Incentives are risk indicators, not evidence that a particular manager committed fraud \citep[chap.~11, solutions~13--20, PDF pp.~218--219]{statement-analysis-workbook-2020}.
\qffield{Practice} Meridian's bonus threshold is based on EBITDA, and its lender tests debt/EBITDA. Capitalizing 20 of routine maintenance improves both the bonus metric and apparent covenant headroom. An independent audit committee should challenge the capitalization criteria, inspect work orders and ask whether spending creates an asset. The auditor's reasonable-assurance work, internal controls, board oversight, securities regulation, lender verification and potential litigation provide different forms of discipline. None guarantees error-free reports: estimates can be uncertain, controls overridden, sampling incomplete and regulators reactive.
\qffield{Conclusion} For the lender, require monthly borrowing-base data and evidence of receivable collection, not only management's adjusted earnings. For the analyst, reconcile incentive metrics to audited figures and compare estimate changes with operations. Strong governance raises confidence in the process but cannot substitute for transaction evidence or make an economically weak business strong. Audit responsibilities include risks of material misstatement from fraud; describing an audit as sampling does not remove that responsibility.
\end{qfmethod}

\begin{qfmethod}[Recognition, estimates, real actions and non-GAAP presentation]\label{meth:earnings-management-choices}
\qffield{Statement} Reported outcomes can be shifted through revenue cut-off, returns/credit-loss estimates, useful lives, impairment timing, capitalization, provisions, acquisition allocation and classification. Real actions such as postponing maintenance or research change the business itself; an accounting estimate changes recognition of that business. A reserve release may boost income without new sales or cash, while a large current provision can depress profit and improve later periods. Compare both the choice and its stated evidence \citep[chap.~11, solutions~8, 10--12 and~22--27, PDF pp.~218--220]{statement-analysis-workbook-2020}.
\qffield{Practice} In the Meridian variant, a 60 sale before control transfers recognizes revenue too early and removes inventory costing 36. Capitalizing 20 of routine maintenance understates operating costs. At a 25 percent rate and with tax recognition deferred as specified in the case, reported net income rises 33 without better economics. Separately, extending remaining asset lives reduces current depreciation but cannot be justified merely by an earnings target. Genuine operational deferral of maintenance might raise current cash flow while damaging future reliability.
\qffield{Presentation} US SEC staff guidance warns that excluding normal recurring cash operating expenses, adjusting charges but not analogous gains, or replacing required revenue recognition with an individually tailored measure can mislead even with a reconciliation. Apply consistent definitions and explain changes across periods \citep[Questions~100.01--100.06, updated 13~December 2022]{non-gaap-disclosures-2022}. Distinguish an analyst's explicitly labeled normalization from a compliant published measure. IFRS 18 adds specified disclosures for management-defined performance measures from its effective period; it does not make every adjusted performance claim a standardized number.
\end{qfmethod}

\begin{qfmethod}[From warning signs to corroborated reporting problems]\label{meth:accounting-warning-sign-investigation}
\qffield{Statement} Screen for receivables growing faster than sales, rising overdue balances, unusual contract assets, inventory accumulation, reserve releases, declining depreciation rates, repeated ``one-time'' adjustments, persistent profit/CFO gaps, supplier-payment delays and related-party transactions. Compare ratios with prior periods, peers, segment mix and business explanations. Statistical screens, including accrual or manipulation scores, identify associations and possible investigation targets, not causation or a legal finding \citep[chap.~11, solutions~23--27, PDF pp.~219--220; chap.~17, solutions~5--9, PDF p.~252]{statement-analysis-workbook-2020}.
\qffield{Practice} In the Meridian reporting variant, revenue is 1,060 and receivables 210, versus prior revenue 900 and receivables 120: growth is 17.78 percent versus 75 percent. That mismatch prompts inspection of delivery terms and subsequent cash receipts. A signed consignment agreement retaining control contradicts management's claim of an earned sale; a customer-confirmed completed delivery would support a legitimate increase instead. Maintenance invoices and asset registers distinguish a genuine improvement from ordinary servicing. Investigate inventory quantities and subsequent prices before interpreting the lower reported inventory of 164 as better efficiency.
\qffield{Method} Preserve a hypothesis and an alternative explanation for each flag, obtain transaction or independent evidence, quantify a proposed correction and test all statement links. If access is insufficient, report the uncertainty and sensitivity rather than manufacturing an adjustment. CFO above net income is one useful observation, not an all-purpose quality test: the combined variant's CFO of 157.5 exceeds its overstated income of 145.5.
\end{qfmethod}

\begin{qfmethod}[Evaluating a financial report with a reconciled correction case]\label{meth:financial-report-quality-case}
\qffield{Statement} A report-quality assessment combines the business model and incentives, the reporting basis, accounting choices, internal consistency, external evidence and the consequences for the intended decision. Evaluate recognition and completeness before trusting ratios. Then separate warranted adjustments from uncertainty that cannot be resolved. The workbook's quality framework emphasizes how biased acquisition values, accruals and presentation can obscure otherwise compliant-looking results \citep[chap.~17, solutions~1--4 and~10--11, PDF pp.~251--252]{statement-analysis-workbook-2020}.
\qffield{Practice} Consider a separate distorted version of Meridian's 2026 accounts. It prematurely recognizes a 60 credit sale and removes stock costing 36; the transaction is not taxable until genuinely earned. It also capitalizes 20 of ordinary maintenance at year-end, with no depreciation yet, although tax deducts that cost immediately. These stipulated facts fail the recognition criteria; no assumption about intent is required. A 25 percent tax rate produces DTL 11. The resulting figures reconcile as follows:
\begin{center}\begin{tabular}{lrr|lrr}
Income/cash flow & Reported & Corrected & Balance sheet & Reported & Corrected\\\hline
Revenue & 1,060 & 1,000 & Cash & 87.5 & 87.5\\
COGS & 636 & 600 & Receivables & 210 & 150\\
Operating expense & 160 & 180 & Inventory & 164 & 200\\
Depreciation & 50 & 50 & PP\&E & 450 & 430\\
Pretax profit & 194 & 150 & Assets & 911.5 & 867.5\\
Tax expense & 48.5 & 37.5 & Other liabilities & 385 & 385\\
Net income & 145.5 & 112.5 & Deferred tax liability & 11 & 0\\
CFO & 157.5 & 137.5 & Equity & 515.5 & 482.5\\
Investing cash flow & $-100$ & $-80$ & Financing cash flow & $-70$ & $-70$
\end{tabular}\end{center}
\qffield{Conclusion} Remove receivables 60, restore inventory 36, remove PP\&E 20 and DTL 11, and reduce retained earnings 33. Income falls 44 before tax and 33 after tax; CFO falls 20 and investing outflow falls 20, leaving total cash unchanged. Both the distorted and corrected balance sheets balance: arithmetic consistency alone did not detect the bad accounting. For a lender, use the corrected operating profit, collateral and cash classifications and seek remediation of the controls; do not issue an opinion on fraud from these hypothetical findings alone.
\end{qfmethod}

\begin{qfmethod}[Earnings persistence, accruals and mean reversion]\label{meth:earnings-persistence-accruals}
\qffield{Statement} Sustainable earnings arise from recurring economics, after the recurring investment needed to support them. Decompose earnings into cash and accrual components, inspect accrual sources, remove genuine transitory items symmetrically and compare margins with competitive conditions. Accruals include legitimate credit sales and matching adjustments as well as estimation errors; a high accrual measure is not by itself proof of manipulation. The workbook reports lower persistence of accrual-heavy earnings as a general empirical tendency, not a law for every firm \citep[chap.~17, solutions~3, 7, 10 and~12, PDF pp.~251--253]{statement-analysis-workbook-2020}.
\qffield{Practice} Meridian's simple cash-flow accrual measure $(NI-CFO)/average\ assets$ is $(112.5-137.5)/833.75=-3.00\%$. It is explicitly a net-income/CFO diagnostic, not the broader net-operating-assets accrual definition that also adjusts investing flows. The negative number partly reflects depreciation and does not establish perpetual cash quality. In the distorted variant it is $(145.5-157.5)/855.75=-1.40\%$: still negative despite recognition errors. The premature sale creates a current after-tax profit increment of 18; when the genuine sale later occurs, that amount cannot be recognized a second time. The resulting reversal illustrates why a temporary accrual boost can cause faster mean reversion than recurring customer cash earnings.
\qffield{Evaluation} Use corrected income 112.5 as the starting point, then challenge the improved operating margin, collection performance and reinvestment needs. Forecast recurring earnings separately from disposal gains, tax releases and restructuring that is actually repeated. A stable cash component may support persistence, but customer concentration, price competition or deferred maintenance can still weaken it. No empirical coefficients, trading profitability or fraud probability are inferred from this synthetic case or from an unread underlying paper.
\end{qfmethod}

\begin{qfmethod}[Cash-flow sustainability, asset recovery and obligation completeness]\label{meth:cash-balance-quality-risk}
\qffield{Statement} Assess cash generation by source, repeatability, timing and classification, then assess whether assets are recoverable and obligations complete. Useful evidence includes cash restrictions, receivable aging, inventory obsolescence, impairment assumptions, debt maturities/covenants, lease commitments, pension funding, guarantees, contingencies and subsequent events. Management commentary, audit reports, credit agreements and regulator correspondence supplement statement totals \citep[chap.~17, solutions~13--19, PDF p.~253; chap.~1, solutions~9--19, PDF p.~174]{statement-analysis-workbook-2020}.
\qffield{Practice} Corrected Meridian CFO 137.5 covers capex 80 but not capex plus principal 30 and dividends 40, so 12.5 comes from cash reserves. That is affordable for one year, but not a sustainable perpetual funding source. An additional 20 received from a receivable transaction that fails derecognition and is a secured borrowing must be financing cash, with the receivable retained and debt recognized. If management places it in CFO, reclassify 20 from operating to financing without changing total cash. A genuine qualifying receivable sale is different and requires examination of the transfer terms.
\qffield{Evaluation} In the distorted variant, cash is real but part of the purported operating cash improvement is maintenance misclassification. Correct both the income and cash-flow analyses. Inventory and receivables must support realizable collateral, while hidden guarantees or near-term debt acceleration may make reported equity a weak buffer. Do not simply net restricted cash against all debt. For the credit decision, require reliable aged receivables/inventory, test payment dates and model a collection slowdown; state the unresolved range where recovery evidence is absent.
\end{qfmethod}
